% !PRINTED 430
\weberpage{430}
Wir müssen uns hier auf eine kurze Angabe der Resultate beschränken, und werden in der Folge keine Anwendungen davon machen, so dass diese Ausführungen, die nur einem mit der Theorie der transcendenten Functionen einigermaassen vertrauten Leser ganz verständlich sein können, ohne Störung des Zusammenhanges auch übergangen werden können.

Es möge $\omega$ eine Variable bedeuten, deren imaginärer Theil positiv ist, so dass
\weq{1}{q = e^{\pi i \omega}}
eine Grösse ist, deren absoluter Werth ein echter Bruch ist. Wir wollen die folgenden drei Functionen dieser Variablen $\omega$ betrachten:
\weq{2}{\eta(\omega) = q^{1/12}\prod_{1,\infty}^{\nu}(1 - q^{2\nu}) = q^{1/12}\sum_{-\infty,\infty}^{\nu}(-1)^\nu q^{3\nu^2+\nu} = \sum_{-\infty,\infty}^{\nu}(-1)^\nu q^{\frac{(6\nu+1)^2}{12}},}
\weq{3}{f(\omega) = q^{-\frac{1}{24}}\prod(1 + q^{2\nu-1}),}
\weq{4}{\gamma_2(\omega) = \frac{f(\omega)^{24} - 16}{f(\omega)^8}\footnote{Man vergl.: \gesperrt{Weber}, \glqq{}Elliptische Functionen und algebraische Zahlen\grqq{}. Braunschweig 1891. S. 63, 86, 149.}.}
In der Transformationstheorie der elliptischen Functionen\footnote{Ebend., S. 263, 316.} wird eine Gleichung 6\textsuperscript{ten} Grades abgeleitet, deren Wurzeln die sechs Grössen
\weq{5}{v_\infty = 5\left(\frac{\eta(5\omega)}{\eta(\omega)}\right)^2, \qquad v_\xi = \left(\frac{\eta\left(\dfrac{-24\xi+\omega}{5}\right)}{\eta(\omega)}\right)^2}
sind, wenn $\xi$ ein volles Restsystem nach dem Modul 5 durchläuft. Diese Gleichung lautet
\weq{6}{v^6 + 10v^3 - \gamma_2(\omega)v + 5 = 0,}
und hat, wie man sieht, grosse Aehnlichkeit mit der im vorigen Paragraphen betrachteten Resolvente 6\textsuperscript{ten} Grades
\weq{7}{U^6 - 10zU^3 + z^2U + 5z^2 = 0.}
Diese beiden Gleichungen gehen geradezu in einander über, wenn man
\weq{8}{z = \gamma_2(\omega)^3}
\weq{9}{U = -\gamma_2(\omega)v}
%CONT>>
